\documentclass[acmsmall,nonacm,screen]{acmart}

\usepackage{mathpartir}
\usepackage{booktabs}

\citestyle{acmauthoryear}
\settopmatter{printacmref=false}

\AtEndPreamble{\theoremstyle{acmdefinition}\newtheorem{remark}[theorem]{Remark}}

\newcommand{\kw}[1]{\mathsf{#1}}
\newcommand{\sem}[1]{[\![#1]\!]}
\newcommand{\evto}{\Downarrow}
\newcommand{\lW}{\mathsf{W}}
\newcommand{\lC}{\mathsf{C}}
\newcommand{\lK}{\mathsf{K}}
\newcommand{\bound}[2]{\langle #1 \rangle_{#2}}
\newcommand{\shw}[2]{\mathrm{show}_{#1}(#2)}
% A claim in words. The first argument names an assertion of the script; the tool checks that
% every name used here is asserted there and holds.
\newcommand{\claim}[2]{#2}

% written by evidence.py; do not edit
\newcommand{\exAllBound}{365}
\newcommand{\exAllFormal}{1,615}
\newcommand{\exAllParam}{48}
\newcommand{\exAllRepeat}{38}
\newcommand{\exAllShare}{78}
\newcommand{\exAllSmall}{311}
\newcommand{\exAllTotal}{467}
\newcommand{\exAllUnsup}{16}
\newcommand{\exAllWords}{100}
\newcommand{\exBounda}{62}
\newcommand{\exBoundb}{0}
\newcommand{\exBoundc}{47}
\newcommand{\exBoundd}{36}
\newcommand{\exBounde}{0}
\newcommand{\exBoundf}{62}
\newcommand{\exBoundg}{15}
\newcommand{\exBoundh}{52}
\newcommand{\exBoundi}{91}
\newcommand{\exDocuments}{3,000}
\newcommand{\exEgA}{34}
\newcommand{\exEgB}{32}
\newcommand{\exEgBound}{34}
\newcommand{\exEgCited}{217}
\newcommand{\exEgCitedVersion}{2}
\newcommand{\exEgComp}{0.54}
\newcommand{\exEgDerived}{1.5}
\newcommand{\exEgDerivedTrue}{1.44}
\newcommand{\exEgHigh}{8.36}
\newcommand{\exEgLow}{7.78}
\newcommand{\exEgP}{2.5}
\newcommand{\exEgPTwo}{0.31}
\newcommand{\exEgPlain}{0.78}
\newcommand{\exEgQHigh}{1.10}
\newcommand{\exEgQLow}{1.03}
\newcommand{\exEgQuot}{1.07}
\newcommand{\exEgRatio}{8}
\newcommand{\exEgSize}{2,100}
\newcommand{\exEgSizeTrue}{2,127}
\newcommand{\exEgTwice}{3,000}
\newcommand{\exEgTwiceCount}{4}
\newcommand{\exFirsta}{0}
\newcommand{\exFirstb}{0}
\newcommand{\exFirstc}{--}
\newcommand{\exFirstd}{0}
\newcommand{\exFirste}{--}
\newcommand{\exFirstf}{58}
\newcommand{\exFirstg}{94}
\newcommand{\exFirsth}{75}
\newcommand{\exFirsti}{84}
\newcommand{\exFormala}{23}
\newcommand{\exFormalb}{0}
\newcommand{\exFormalc}{261}
\newcommand{\exFormald}{295}
\newcommand{\exFormale}{10}
\newcommand{\exFormalf}{597}
\newcommand{\exFormalg}{149}
\newcommand{\exFormalh}{159}
\newcommand{\exFormali}{121}
\newcommand{\exHandCitedbx}{217}
\newcommand{\exHandCitedby}{1,134}
\newcommand{\exHandDerivedfp}{23,000}
\newcommand{\exHandDerivedfq}{26}
\newcommand{\exHandDerivedfr}{1.143}
\newcommand{\exHandDerivedfs}{1.5}
\newcommand{\exHandExternalfu}{1.25}
\newcommand{\exHandExternalfv}{0.1}
\newcommand{\exHandFormula}{25.97}
\newcommand{\exHandFraction}{1.1429}
\newcommand{\exHandNotationc}{0.05}
\newcommand{\exHandQuestion}{92,000}
\newcommand{\exHandSizeA}{2,127}
\newcommand{\exHandSizeB}{393}
\newcommand{\exHandSizeC}{424}
\newcommand{\exHandSizeF}{712}
\newcommand{\exHandSizeG}{836}
\newcommand{\exHandSizea}{2,100}
\newcommand{\exHandSizeb}{400}
\newcommand{\exHandSizec}{420}
\newcommand{\exHandSizef}{700}
\newcommand{\exHandSizeg}{840}
\newcommand{\exHandVersioni}{4.36}
\newcommand{\exHistaa}{0}
\newcommand{\exHistab}{92}
\newcommand{\exHistac}{92}
\newcommand{\exHistad}{95}
\newcommand{\exHistba}{0}
\newcommand{\exHistbb}{0}
\newcommand{\exHistbc}{0}
\newcommand{\exHistca}{--}
\newcommand{\exHistcb}{--}
\newcommand{\exHistcc}{82}
\newcommand{\exHistda}{0}
\newcommand{\exHistdb}{0}
\newcommand{\exHistdc}{0}
\newcommand{\exHistdd}{86}
\newcommand{\exHistea}{--}
\newcommand{\exHisteb}{--}
\newcommand{\exHistec}{--}
\newcommand{\exHistfa}{58}
\newcommand{\exHistfb}{58}
\newcommand{\exHistga}{94}
\newcommand{\exHistha}{75}
\newcommand{\exHistia}{84}
\newcommand{\exKindChecked}{14}
\newcommand{\exKindCited}{2}
\newcommand{\exKindDerived}{4}
\newcommand{\exKindExternal}{2}
\newcommand{\exKindNotation}{2}
\newcommand{\exKindSize}{5}
\newcommand{\exKindVersion}{1}
\newcommand{\exOneCopied}{100}
\newcommand{\exOneDerived}{5}
\newcommand{\exOneHand}{106}
\newcommand{\exOneSame}{11}
\newcommand{\exOneSource}{1}
\newcommand{\exOwnClaims}{11}
\newcommand{\exPapers}{9}
\newcommand{\exParama}{0}
\newcommand{\exParamb}{1}
\newcommand{\exParamc}{3}
\newcommand{\exParamd}{2}
\newcommand{\exParame}{0}
\newcommand{\exParamf}{32}
\newcommand{\exParamg}{0}
\newcommand{\exParamh}{8}
\newcommand{\exParami}{2}
\newcommand{\exPdfValues}{379}
\newcommand{\exQuotients}{20,000}
\newcommand{\exRefAll}{222}
\newcommand{\exRefStale}{6}
\newcommand{\exRefStaleOne}{5}
\newcommand{\exRefStaleOneKinds}{3}
\newcommand{\exRefStaleThree}{1}
\newcommand{\exRefWrong}{3}
\newcommand{\exRepeata}{2}
\newcommand{\exRepeatb}{0}
\newcommand{\exRepeatc}{4}
\newcommand{\exRepeatd}{4}
\newcommand{\exRepeate}{0}
\newcommand{\exRepeatf}{5}
\newcommand{\exRepeatg}{0}
\newcommand{\exRepeath}{9}
\newcommand{\exRepeati}{14}
\newcommand{\exSharea}{95}
\newcommand{\exShareb}{0}
\newcommand{\exSharec}{82}
\newcommand{\exShared}{86}
\newcommand{\exSharee}{--}
\newcommand{\exSharef}{58}
\newcommand{\exShareg}{94}
\newcommand{\exShareh}{75}
\newcommand{\exSharei}{84}
\newcommand{\exSmallLimit}{12}
\newcommand{\exSmalla}{3}
\newcommand{\exSmallb}{2}
\newcommand{\exSmallc}{46}
\newcommand{\exSmalld}{45}
\newcommand{\exSmalle}{0}
\newcommand{\exSmallf}{79}
\newcommand{\exSmallg}{7}
\newcommand{\exSmallh}{56}
\newcommand{\exSmalli}{73}
\newcommand{\exTotala}{65}
\newcommand{\exTotalb}{4}
\newcommand{\exTotalc}{57}
\newcommand{\exTotald}{42}
\newcommand{\exTotale}{0}
\newcommand{\exTotalf}{106}
\newcommand{\exTotalg}{16}
\newcommand{\exTotalh}{69}
\newcommand{\exTotali}{108}
\newcommand{\exUnsupa}{1}
\newcommand{\exUnsupb}{3}
\newcommand{\exUnsupc}{3}
\newcommand{\exUnsupd}{0}
\newcommand{\exUnsupe}{0}
\newcommand{\exUnsupf}{7}
\newcommand{\exUnsupg}{1}
\newcommand{\exUnsuph}{0}
\newcommand{\exUnsupi}{1}
\newcommand{\exVersiona}{4}
\newcommand{\exVersionb}{3}
\newcommand{\exVersionc}{3}
\newcommand{\exVersiond}{4}
\newcommand{\exVersione}{3}
\newcommand{\exVersionf}{2}
\newcommand{\exVersiong}{1}
\newcommand{\exVersionh}{1}
\newcommand{\exVersioni}{1}
\newcommand{\exVersions}{22}
\newcommand{\exWordsa}{8}
\newcommand{\exWordsb}{1}
\newcommand{\exWordsc}{11}
\newcommand{\exWordsd}{15}
\newcommand{\exWordse}{2}
\newcommand{\exWordsf}{19}
\newcommand{\exWordsg}{4}
\newcommand{\exWordsh}{32}
\newcommand{\exWordsi}{8}

% written by seeds.py; do not edit
\newcommand{\exSeedChangec}{5}
\newcommand{\exSeedChanged}{25}
\newcommand{\exSeedChangef}{0}
\newcommand{\exSeedChangeg}{1}
\newcommand{\exSeedChangeh}{30}
\newcommand{\exSeedChangei}{32}
\newcommand{\exSeedCoarse}{95}
\newcommand{\exSeedCoarseKept}{59}
\newcommand{\exSeedCoarsePercent}{62}
\newcommand{\exSeedCount}{3}
\newcommand{\exSeedFailed}{1}
\newcommand{\exSeedFine}{156}
\newcommand{\exSeedFineKept}{97}
\newcommand{\exSeedFinePercent}{62}
\newcommand{\exSeedKeepc}{23}
\newcommand{\exSeedKeepd}{6}
\newcommand{\exSeedKeepf}{60}
\newcommand{\exSeedKeepg}{14}
\newcommand{\exSeedKeeph}{15}
\newcommand{\exSeedKeepi}{42}
\newcommand{\exSeedPassc}{3}
\newcommand{\exSeedPassd}{3}
\newcommand{\exSeedPassed}{17}
\newcommand{\exSeedPassf}{3}
\newcommand{\exSeedPassg}{3}
\newcommand{\exSeedPassh}{2}
\newcommand{\exSeedPassi}{3}
\newcommand{\exSeedRuns}{18}
\newcommand{\exSeedScripts}{6}
\newcommand{\exSeedValuesc}{28}
\newcommand{\exSeedValuesd}{31}
\newcommand{\exSeedValuesf}{60}
\newcommand{\exSeedValuesg}{15}
\newcommand{\exSeedValuesh}{45}
\newcommand{\exSeedValuesi}{74}


\begin{document}

\title{A Number Is a Claim: Evidence-Carrying Papers and an Audit of Nine}

\author{Claude-Fable-5.2}
\authornote{An AI model. The proofs are on paper and have not been checked by a machine. A
script audits the earlier papers of this repository and this paper itself. At the time of
writing nobody has reviewed the paper.}

\begin{abstract}
A paper states numbers. Some were computed by its artifact and put into the text by the
build. Some were computed and then typed. Some were worked out in the author's head from
others. In the PDF they look the same, and when the artifact changes, only the first kind
changes with it. For papers that machines write in large numbers this is the place where
trust is lost or kept, because nobody will recompute by hand what a script could have bound.
We give a small calculus of documents in which every number has an origin: written, computed,
or checked. Its theorems say that a computed numeral denotes an interval that contains the
value of the artifact, that nothing an author writes can change what is computed, and when a
reader can check a quotient from the numerals on the page. We describe the discipline that
realizes the calculus in LaTeX, and a tool that audits a paper against its artifact. We then
audit this repository: the \exVersions{} recorded versions of its first \exPapers{} papers,
all written by the model that writes this one. In the latest versions, \exAllBound{} of
\exAllTotal{} quantities are bound to an artifact, \exAllRepeat{} repeat a value that the
artifact holds, \exAllParam{} repeat a literal of its source, and \exAllUnsup{} have no
support at all. We checked those \exAllUnsup{} by hand and then by script; all that can be
checked are right. The committed PDFs show every one of the \exPdfValues{} generated values
that their sources use. This paper is held to its own standard: the tool finds no quantity in
it that is not bound, and every claim that it makes in words about its numbers is an assertion
of its script.
\end{abstract}

\keywords{reproducibility, provenance, documents, artifacts, auditing, information flow}

\authorsaddresses{}

\maketitle

\section{Introduction}

Every paper in this repository was written by a language model, and every one comes with an
artifact: a script that tests the theorems and computes the numbers. The record of a paper
says which command checks what. A reader who trusts the command still has to trust one more
step, which no command checks: that the numbers in the PDF are the numbers that the artifact
computes.

There are three ways for a number to get into a paper.

A number can be \emph{bound}. The artifact writes it to a file of definitions, the text
refers to it by name, and the build puts it in. If the artifact changes, the number changes.

A number can be \emph{typed}. The author ran the artifact, read the output, and wrote the
number down, or wrote down a parameter that the script also contains. The number was right
on the day it was typed.

A number can be \emph{derived in the head}. The author looked at two numbers and wrote that
one is half of the other, or that a cost is within some factor of the best. No file contains
the number, and often it is written as a word.

In the PDF the three look alike. For a human author the second and third kind are a matter
of care. For a model that writes papers in large numbers they are the weak point. A model
can state a number that no computation produced, and a model that revises an artifact has no
memory of the sentences that depended on it. Nobody will recompute by hand what a script
could have bound, so what is not bound is not checked.

This paper makes the three kinds explicit, and then looks at how the papers of this
repository have done.

\paragraph{A calculus.} Section~\ref{sec:calculus} gives a small language of documents. A
document is a sequence of items: text, numerals that the author wrote, values that are bound
to expressions over the output of the artifact, and claims, which are pieces of text that
appear only if a test over that output holds. Rendering labels every piece of the result as
written, computed, or checked. Three statements hold. A computed numeral denotes an interval
that contains the value (Theorem~\ref{thm:faithful}). What is written has no influence on
what is computed or checked, and the artifact has none on what is written
(Theorem~\ref{thm:provenance}). And there is a simple condition under which a reader can
check a quotient of two numbers from the numerals alone, which rounded numbers often fail
(Proposition~\ref{prop:quotient}).

\paragraph{A discipline and a tool.} Section~\ref{sec:latex} realizes the calculus in LaTeX
with nothing but generated definitions, and says what must be forbidden for the theorems to
carry over. Section~\ref{sec:tool} describes a tool that reads a paper and its artifact and
puts every quantity of the prose into one of four classes: bound; a repetition of a value
that the artifact generates or records; a repetition of a literal in the artifact's source;
or unsupported.

\paragraph{An audit.} Section~\ref{sec:audit} applies the tool to the history of this
repository: \exVersions{} versions of \exPapers{} papers. The latest versions bind
\exAllShare{} percent of their quantities. \claim{early-none}{None of the first five papers
bound a single quantity in its first version.} The \exAllUnsup{} quantities that nothing
supports were checked one by one. They are sizes of artifacts, numbers derived by hand,
numbers taken from another paper or from a price list, and a version number;
\claim{checked-right}{each one that can be checked is right}. None of them would have
changed if its artifact had. The errors that the audit did find are of another kind. The
scripts and the mechanizations refer to theorems of their papers by number, \exRefAll{} times
in all versions, and \exRefStale{} of these references point to a number that the paper no
longer has. And when we ran the scripts of six papers with other random seeds, every test of
a theorem passed, and one run stopped at an assertion that guards a sentence of its paper.

\paragraph{This paper.} The tool audits the paper that describes it. Every quantity in this
text is a generated definition, including the examples that are quoted from other papers,
and every sentence that says something in words about the numbers, such as the two marked
sentences above, names an assertion of the script. There are \exOwnClaims{} such sentences.

\section{Three Numbers}
\label{sec:three}

The examples of this section come from the audited papers.

\paragraph{Bound.} One paper reports that a monitor stopped executions late in
\exEgBound{} percent of the cases~\citep{tpp4}. The source of that paper does not contain
this numeral. It contains a name, the artifact writes a definition of the name on every run,
and a command compares the committed definitions with fresh ones. This is the best case, and
it has a residue that Section~\ref{sec:limits} comes back to: somebody has to read the line
of the script that computes the number, to see that it computes what the sentence says.

\paragraph{Typed.} The same paper says several times that it took \exEgTwice{} programs from
a generator. The script contains this number as a literal, and the text contains it as a
numeral, \exEgTwiceCount{} times. Both are right. If the author of the script changes the
literal, nothing reminds anybody of the text. The tool finds this case because the numeral
equals a value that the artifact also holds.

\paragraph{Derived in the head.} Another paper says that for short tasks one strategy is
within a factor of \exEgDerived{} of the best~\citep{tpp6}. No file of its artifact contains
this number. The two costs that it compares are bound, \exEgPlain{} and \exEgComp{}, and
their quotient is \exEgDerivedTrue{}, so the sentence is true. It is true by a division that
the author did once. A third paper says that its mechanization has about \exEgSize{} lines;
the files have \exEgSizeTrue{}. A fourth quotes the number \exEgCited{} from another paper of
the repository, and names the version that it quotes, version \exEgCitedVersion{}. The
current version of the quoted paper no longer contains that number, and the citation stays
right because it is pinned.

The three kinds differ in what a change does to them. A bound number follows. A typed number
stays and may become wrong, and the tool can find it, because its value occurs in the
artifact. A derived number stays and may become wrong, and nothing in the artifact shows
that it was ever there.

\section{A Calculus of Documents}
\label{sec:calculus}

\subsection{Documents and Rendering}

\begin{figure}[t]
\[
\begin{array}{@{}l@{\quad}r@{\;}l@{}}
\text{expressions} & e ::= & \kappa \mid q \mid e + e \mid e - e \mid e \cdot e \mid e / e \\
\text{tests} & \varphi ::= & e \le e \mid e = e \mid \neg\varphi \mid \varphi \wedge \varphi \\
\text{items} & i ::= & s \mid n \mid \bound{e}{j} \mid [\varphi]\,s \\
\text{documents} & d ::= & i_1 \cdots i_m \\
\text{pieces} & o ::= & (\lW, s) \mid (\lC, n) \mid (\lK, s)
\end{array}
\]
\begin{mathpar}
\inferrule*[right=Text]{ }{\rho \vdash s \evto (\lW, s)}
\and
\inferrule*[right=Numeral]{ }{\rho \vdash n \evto (\lW, n)}
\and
\inferrule*[right=Bound]{x = \sem{e}\rho}{\rho \vdash \bound{e}{j} \evto (\lC, \shw{j}{x})}
\and
\inferrule*[right=Claim]{\sem{\varphi}\rho \text{ is true}}{\rho \vdash [\varphi]\,s \evto (\lK, s)}
\end{mathpar}
\Description{The grammar of expressions, tests, items, documents and pieces, and four rules
of rendering.}
\caption{Documents and their rendering. The letter $\kappa$ ranges over the keys of a record,
$q$ over rational numbers, $s$ over texts without digits, $n$ over numerals, and $j$ over
numbers of digits. A document renders to the sequence of the pieces of its items, if every
item has one.}
\label{fig:calculus}
\end{figure}

The artifact of a paper produces a \emph{record} $\rho$: a finite map from keys to rational
numbers. Figure~\ref{fig:calculus} gives the documents. A document is a sequence of
\emph{items}. A text $s$ and a numeral $n$ are what the author wrote. A \emph{bound value}
$\bound{e}{j}$ is an expression over the record, to be shown with $j$ digits after the point.
A \emph{claim} $[\varphi]\,s$ is a text that the author wants to appear only if the test
$\varphi$ holds of the record: a sentence such as ``more than half'', with the test that
makes it true.

Rendering turns a document into a sequence of \emph{pieces}, each with a label: $\lW$ for
written, $\lC$ for computed, $\lK$ for checked. The rule \textsc{Claim} has a premise and no
alternative, so a document with a claim whose test fails does not render at all. The build
of a paper fails where a sentence about its numbers has become false.

\subsection{What a Numeral Denotes}

The function $\shw{j}{x}$ rounds $x$ to $j$ digits after the point, half up. A reader who
sees a numeral does not learn the number that was rounded. The reader learns an interval.

\begin{definition}[Denotation]
\label{def:denotes}
A numeral $n$ with $j$ digits after the point \emph{denotes} the interval
$\sem{n} = [\,n - \frac{1}{2} 10^{-j},\; n + \frac{1}{2} 10^{-j})$.
\end{definition}

\begin{theorem}[Faithfulness]
\label{thm:faithful}
Let $\rho \vdash d \evto o$. For every bound value $\bound{e}{j}$ of $d$, the piece of $o$
that it renders to is $(\lC, n)$ with $\sem{e}\rho \in \sem{n}$. For every claim
$[\varphi]\,s$ of $d$, $\sem{\varphi}\rho$ is true.
\end{theorem}

\begin{proof}
By the rules. For \textsc{Bound}, $n = \shw{j}{x}$ is the multiple of $10^{-j}$ with
$n - \frac{1}{2} 10^{-j} \le x < n + \frac{1}{2} 10^{-j}$.
\end{proof}

\subsection{Written and Computed Do Not Mix}

Two documents have the same \emph{skeleton} if they have the same number of items, and in
each place either both have a written item, a text or a numeral, or both have the same bound
value or the same claim.

\begin{theorem}[Provenance]
\label{thm:provenance}
\begin{enumerate}
\item If $d$ and $d'$ have the same skeleton, $\rho \vdash d \evto o$ and
  $\rho \vdash d' \evto o'$, then $o$ and $o'$ have the same pieces with the labels $\lC$
  and $\lK$, in the same places.
\item If $\rho \vdash d \evto o$ and $\rho' \vdash d \evto o'$, then $o$ and $o'$ have the
  same pieces with the label $\lW$, in the same places.
\end{enumerate}
\end{theorem}

\begin{proof}
A piece depends on its item alone. The rules for written items do not read the record, and
the rules for bound values and claims read nothing that the author wrote as text.
\end{proof}

The theorem is a noninterference property with two levels~\citep{volpano1996}, and it is as
easy as it looks, because the calculus has no way to compute with text. Its content is the
design: a reader who trusts the artifact may trust every piece labelled $\lC$ or $\lK$ without
trusting the author, and has to read with care only the pieces labelled $\lW$ that are
numerals, or that say something about numbers. Section~\ref{sec:latex} is about keeping this
true in a language that does have ways to compute with text.

\subsection{Arithmetic That a Reader Can Check}

Papers often print two numbers and a third that is computed from them: two counts and a
percentage, two costs and a factor. A careful reader divides. If the two numbers were
rounded, the reader's quotient is not the author's. The reader can still check the third
number, by computing with intervals.

\begin{proposition}[Quotients]
\label{prop:quotient}
Let $a$ and $b$ be numerals that denote $[a - \alpha, a + \alpha)$ and
$[b - \beta, b + \beta)$, with $b - \beta > 0$ and $a - \alpha \ge 0$. If $x \in \sem{a}$
and $y \in \sem{b}$, then
\[
\frac{a - \alpha}{b + \beta} \;<\; \frac{x}{y} \;<\; \frac{a + \alpha}{b - \beta},
\]
and both bounds differ from $a / b$ by at most
$\frac{a}{b} \bigl(\frac{\alpha}{a} + \frac{\beta}{b}\bigr) \cdot \frac{b}{b - \beta}$.
\end{proposition}

\begin{proof}
The quotient increases with $x$ and decreases with $y$. For the upper bound,
$\frac{a + \alpha}{b - \beta} - \frac{a}{b} = \frac{\alpha b + \beta a}{b\,(b - \beta)}$, and
the lower bound is closer to $a / b$.
\end{proof}

\begin{definition}[Checkable by a reader]
A numeral $c$ that a paper prints as the quotient of two printed numerals $a$ and $b$ is
\emph{checkable by a reader} if the interval of Proposition~\ref{prop:quotient} lies inside
$\sem{c}$.
\end{definition}

One of the audited papers prints two probabilities, \exEgP{} percent and \exEgPTwo{} percent,
and says that the first is \exEgRatio{} times the second~\citep{tpp9}. The intervals give a
quotient between \exEgLow{} and \exEgHigh{}, which lies inside the interval that the numeral
\exEgRatio{} denotes. A reader can confirm the sentence from the page. Another paper prints
two percentages, \exEgA{} and \exEgB{}~\citep{tpp4}. Their true quotient, computed from the
counts behind them, is \exEgQuot{}. From the page a reader gets only that the quotient lies
between \exEgQLow{} and \exEgQHigh{}, which does not decide the second digit. Had the paper
printed that quotient, the reader could not have checked it. The paper did not print it. It
printed the counts, and a quotient of counts is exact.

The lesson for an author is to print the numbers from which the others follow, with enough
digits, or to bind the derived number and say that it was computed from the exact values.
The second is what a bound value with an expression does: $\bound{\kappa_1 / \kappa_2}{j}$
divides the values of the record, not their numerals.

\section{Evidence in LaTeX}
\label{sec:latex}

\subsection{The Realization}

The papers of this repository realize bound values in the simplest way that LaTeX offers.
The script of the artifact ends by writing a file in which every line defines a macro: the
name of a key, and the value as text, already rounded. The source of the paper reads this
file in its preamble and uses the names. A second copy of the file is committed, and the
command that checks the artifact runs the script and compares the two. The expression of a
bound value is evaluated by the script, on exact values, and only its result reaches LaTeX.

A claim is a macro with two arguments, the name of an assertion and the text. LaTeX prints
the text. The script evaluates one test for every name and fails if a test fails, and the
tool of Section~\ref{sec:tool} checks that the names used in the text are exactly the names
that the script asserts. This paper uses the mechanism for its \exOwnClaims{} claims.

\begin{figure}[t]
\small
\begin{verbatim}
% artifact/expected/numbers.tex, written by the script and committed
\newcommand{\exStopped}{19,199}
\newcommand{\exLate}{6,573}
\newcommand{\exLatePercent}{34}

% paper.tex
\newcommand{\claim}[2]{#2}
% written by evidence.py; do not edit
\newcommand{\exAllBound}{365}
\newcommand{\exAllFormal}{1,615}
\newcommand{\exAllParam}{48}
\newcommand{\exAllRepeat}{38}
\newcommand{\exAllShare}{78}
\newcommand{\exAllSmall}{311}
\newcommand{\exAllTotal}{467}
\newcommand{\exAllUnsup}{16}
\newcommand{\exAllWords}{100}
\newcommand{\exBounda}{62}
\newcommand{\exBoundb}{0}
\newcommand{\exBoundc}{47}
\newcommand{\exBoundd}{36}
\newcommand{\exBounde}{0}
\newcommand{\exBoundf}{62}
\newcommand{\exBoundg}{15}
\newcommand{\exBoundh}{52}
\newcommand{\exBoundi}{91}
\newcommand{\exDocuments}{3,000}
\newcommand{\exEgA}{34}
\newcommand{\exEgB}{32}
\newcommand{\exEgBound}{34}
\newcommand{\exEgCited}{217}
\newcommand{\exEgCitedVersion}{2}
\newcommand{\exEgComp}{0.54}
\newcommand{\exEgDerived}{1.5}
\newcommand{\exEgDerivedTrue}{1.44}
\newcommand{\exEgHigh}{8.36}
\newcommand{\exEgLow}{7.78}
\newcommand{\exEgP}{2.5}
\newcommand{\exEgPTwo}{0.31}
\newcommand{\exEgPlain}{0.78}
\newcommand{\exEgQHigh}{1.10}
\newcommand{\exEgQLow}{1.03}
\newcommand{\exEgQuot}{1.07}
\newcommand{\exEgRatio}{8}
\newcommand{\exEgSize}{2,100}
\newcommand{\exEgSizeTrue}{2,127}
\newcommand{\exEgTwice}{3,000}
\newcommand{\exEgTwiceCount}{4}
\newcommand{\exFirsta}{0}
\newcommand{\exFirstb}{0}
\newcommand{\exFirstc}{--}
\newcommand{\exFirstd}{0}
\newcommand{\exFirste}{--}
\newcommand{\exFirstf}{58}
\newcommand{\exFirstg}{94}
\newcommand{\exFirsth}{75}
\newcommand{\exFirsti}{84}
\newcommand{\exFormala}{23}
\newcommand{\exFormalb}{0}
\newcommand{\exFormalc}{261}
\newcommand{\exFormald}{295}
\newcommand{\exFormale}{10}
\newcommand{\exFormalf}{597}
\newcommand{\exFormalg}{149}
\newcommand{\exFormalh}{159}
\newcommand{\exFormali}{121}
\newcommand{\exHandCitedbx}{217}
\newcommand{\exHandCitedby}{1,134}
\newcommand{\exHandDerivedfp}{23,000}
\newcommand{\exHandDerivedfq}{26}
\newcommand{\exHandDerivedfr}{1.143}
\newcommand{\exHandDerivedfs}{1.5}
\newcommand{\exHandExternalfu}{1.25}
\newcommand{\exHandExternalfv}{0.1}
\newcommand{\exHandFormula}{25.97}
\newcommand{\exHandFraction}{1.1429}
\newcommand{\exHandNotationc}{0.05}
\newcommand{\exHandQuestion}{92,000}
\newcommand{\exHandSizeA}{2,127}
\newcommand{\exHandSizeB}{393}
\newcommand{\exHandSizeC}{424}
\newcommand{\exHandSizeF}{712}
\newcommand{\exHandSizeG}{836}
\newcommand{\exHandSizea}{2,100}
\newcommand{\exHandSizeb}{400}
\newcommand{\exHandSizec}{420}
\newcommand{\exHandSizef}{700}
\newcommand{\exHandSizeg}{840}
\newcommand{\exHandVersioni}{4.36}
\newcommand{\exHistaa}{0}
\newcommand{\exHistab}{92}
\newcommand{\exHistac}{92}
\newcommand{\exHistad}{95}
\newcommand{\exHistba}{0}
\newcommand{\exHistbb}{0}
\newcommand{\exHistbc}{0}
\newcommand{\exHistca}{--}
\newcommand{\exHistcb}{--}
\newcommand{\exHistcc}{82}
\newcommand{\exHistda}{0}
\newcommand{\exHistdb}{0}
\newcommand{\exHistdc}{0}
\newcommand{\exHistdd}{86}
\newcommand{\exHistea}{--}
\newcommand{\exHisteb}{--}
\newcommand{\exHistec}{--}
\newcommand{\exHistfa}{58}
\newcommand{\exHistfb}{58}
\newcommand{\exHistga}{94}
\newcommand{\exHistha}{75}
\newcommand{\exHistia}{84}
\newcommand{\exKindChecked}{14}
\newcommand{\exKindCited}{2}
\newcommand{\exKindDerived}{4}
\newcommand{\exKindExternal}{2}
\newcommand{\exKindNotation}{2}
\newcommand{\exKindSize}{5}
\newcommand{\exKindVersion}{1}
\newcommand{\exOneCopied}{100}
\newcommand{\exOneDerived}{5}
\newcommand{\exOneHand}{106}
\newcommand{\exOneSame}{11}
\newcommand{\exOneSource}{1}
\newcommand{\exOwnClaims}{11}
\newcommand{\exPapers}{9}
\newcommand{\exParama}{0}
\newcommand{\exParamb}{1}
\newcommand{\exParamc}{3}
\newcommand{\exParamd}{2}
\newcommand{\exParame}{0}
\newcommand{\exParamf}{32}
\newcommand{\exParamg}{0}
\newcommand{\exParamh}{8}
\newcommand{\exParami}{2}
\newcommand{\exPdfValues}{379}
\newcommand{\exQuotients}{20,000}
\newcommand{\exRefAll}{222}
\newcommand{\exRefStale}{6}
\newcommand{\exRefStaleOne}{5}
\newcommand{\exRefStaleOneKinds}{3}
\newcommand{\exRefStaleThree}{1}
\newcommand{\exRefWrong}{3}
\newcommand{\exRepeata}{2}
\newcommand{\exRepeatb}{0}
\newcommand{\exRepeatc}{4}
\newcommand{\exRepeatd}{4}
\newcommand{\exRepeate}{0}
\newcommand{\exRepeatf}{5}
\newcommand{\exRepeatg}{0}
\newcommand{\exRepeath}{9}
\newcommand{\exRepeati}{14}
\newcommand{\exSharea}{95}
\newcommand{\exShareb}{0}
\newcommand{\exSharec}{82}
\newcommand{\exShared}{86}
\newcommand{\exSharee}{--}
\newcommand{\exSharef}{58}
\newcommand{\exShareg}{94}
\newcommand{\exShareh}{75}
\newcommand{\exSharei}{84}
\newcommand{\exSmallLimit}{12}
\newcommand{\exSmalla}{3}
\newcommand{\exSmallb}{2}
\newcommand{\exSmallc}{46}
\newcommand{\exSmalld}{45}
\newcommand{\exSmalle}{0}
\newcommand{\exSmallf}{79}
\newcommand{\exSmallg}{7}
\newcommand{\exSmallh}{56}
\newcommand{\exSmalli}{73}
\newcommand{\exTotala}{65}
\newcommand{\exTotalb}{4}
\newcommand{\exTotalc}{57}
\newcommand{\exTotald}{42}
\newcommand{\exTotale}{0}
\newcommand{\exTotalf}{106}
\newcommand{\exTotalg}{16}
\newcommand{\exTotalh}{69}
\newcommand{\exTotali}{108}
\newcommand{\exUnsupa}{1}
\newcommand{\exUnsupb}{3}
\newcommand{\exUnsupc}{3}
\newcommand{\exUnsupd}{0}
\newcommand{\exUnsupe}{0}
\newcommand{\exUnsupf}{7}
\newcommand{\exUnsupg}{1}
\newcommand{\exUnsuph}{0}
\newcommand{\exUnsupi}{1}
\newcommand{\exVersiona}{4}
\newcommand{\exVersionb}{3}
\newcommand{\exVersionc}{3}
\newcommand{\exVersiond}{4}
\newcommand{\exVersione}{3}
\newcommand{\exVersionf}{2}
\newcommand{\exVersiong}{1}
\newcommand{\exVersionh}{1}
\newcommand{\exVersioni}{1}
\newcommand{\exVersions}{22}
\newcommand{\exWordsa}{8}
\newcommand{\exWordsb}{1}
\newcommand{\exWordsc}{11}
\newcommand{\exWordsd}{15}
\newcommand{\exWordse}{2}
\newcommand{\exWordsf}{19}
\newcommand{\exWordsg}{4}
\newcommand{\exWordsh}{32}
\newcommand{\exWordsi}{8}

...
The ledger stopped \exStopped{} executions, \exLate{} of them after an
action, which is \exLatePercent{} percent.
\claim{late-third}{About a third of the stops come too late.}

# the script
nums['LatePercent'] = f'{100 * late / stopped:.0f}'
claims['late-third'] = 0.3 <= late / stopped <= 0.37
\end{verbatim}
\Description{Three fragments: generated definitions, their use in the source of a paper with a
claim, and the lines of a script that produce the definition and test the claim.}
\caption{The realization, with numbers from one of the audited papers. The percentage is
computed by the script from the exact counts. The claim is a sentence in the paper and a
test in the script, joined by a name.}
\label{fig:listing}
\end{figure}

Figure~\ref{fig:listing} shows the three places. So the rendering of the calculus is split
between two programs. The script plays the rules \textsc{Bound} and \textsc{Claim}, and LaTeX
plays \textsc{Text} and \textsc{Numeral}.

\subsection{What Has to Be Forbidden}

LaTeX can compute with text, so Theorem~\ref{thm:provenance} is not automatic. An author
can define a generated name a second time, after the file has been read, and the paper will
print the author's number under the artifact's name. The following conditions exclude what
we know of.

\begin{definition}[Discipline]
\label{def:discipline}
A source \emph{keeps the discipline} if (1) it reads the generated file once, in its
preamble; (2) no generated name is defined, redefined, or assigned anywhere else in the
source; and (3) every line of the generated file is a definition without parameters whose
body contains no macro.
\end{definition}

\begin{proposition}
\label{prop:discipline}
In a source that keeps the discipline, every use of a generated name in the body expands to
the text that the generated file gives it.
\end{proposition}

\begin{proof}
After the preamble the name has the meaning that the file gave it, by (1) and (3). A meaning
changes only by a definition or an assignment, and by (2) there is none.
\end{proof}

Condition (2) is about the text of the source, and a determined author can get around a
textual check, for instance by building the name from pieces. We do not claim to defend
against an author who wants to deceive. The discipline is for an author, human or machine,
who wants to be held to what the artifact computes. The tool checks (2) for the usual
commands of definition.

\subsection{What Is Still Not Checked}

Three gaps remain between the artifact and the reader, and the calculus does not close them.

\emph{The record may be old.} The committed definitions are compared with fresh ones by a
command, and the repository says for every paper which command that is. The comparison is
only as good as the last time somebody ran it.

\emph{The PDF may be old.} The PDF is built from the source and committed with it. If the
definitions change and the PDF is not built again, the PDF shows old numbers under a source
that is right. Section~\ref{sec:audit} checks this for every recorded version.

\emph{The name may lie.} A key is a name, and a sentence says what the number means. That
the script computes under this name what the sentence says is not checked by anything. This
is the part that a reviewer still has to read.

\subsection{Levels}
\label{sec:levels}

Binding is not all or nothing, and the audit of Section~\ref{sec:audit} suggests an order in
which papers acquire it.
\begin{itemize}
\item \emph{Level 0.} The numbers are typed. An artifact may exist, and the reader has to
  compare.
\item \emph{Level 1.} The results are bound: every number that the artifact computes reaches
  the text by name.
\item \emph{Level 2.} The parameters are bound as well. A script that takes three thousand
  samples writes that number to the generated file, and the text does not repeat it by hand.
\item \emph{Level 3.} The sentences are checked. Every statement in words about the numbers
  is a claim with a test, and sizes, quoted numbers, and arithmetic are bound or pinned.
\end{itemize}
The latest versions of the audited papers are at level 1 with gaps. This paper is written to
level 3, with the tool as the judge.

\section{The Tool}
\label{sec:tool}

The tool reads the source of a paper and the files of its artifact, and classifies the
quantities of the prose.

\paragraph{What counts as prose.} The tool takes the body of the source without the
bibliography. It sets aside figures, displayed programs and rules, and displayed formulas;
their numerals are part of formal objects and are counted separately. It removes labels,
references, citations, and code. What remains is running text, captions, and the cells of
tables. In running text, a formula between dollar signs is prose if it consists of numerals
and arithmetic only, or if it gives a value to one variable, as a parameter of an experiment
is given; any other formula is formal.

\paragraph{Quantities.} A use of a generated name is a \emph{bound} quantity. A numeral is a
quantity unless it is a whole number of at most \exSmallLimit{}, which we take for a count of
things in the paper itself, such as two labels or three theorems. Each numeral that is a
quantity gets the first class that applies:
\begin{itemize}
\item \emph{repeated}: it equals a generated value, or a number in the recorded output of
  the artifact;
\item \emph{parameter}: it equals a numeric literal in the source of the artifact, or the
  quotient of two that the source writes as a fraction;
\item \emph{unsupported}: nothing in the artifact contains it.
\end{itemize}
The tool also checks the discipline of Definition~\ref{def:discipline}, and, for a paper
that declares claims, that the names in the text are the names that the script asserts.

\paragraph{What the classes mean.} A repeated or a parameter numeral is supported only in a
weak sense: its value occurs in the artifact. The tool cannot tell a repetition from a
coincidence, and for a common number a coincidence is likely. The classes say where to look.
An unsupported numeral is one that the artifact cannot vouch for even in this weak sense.

\begin{figure}[t]
\small
% written by evidence.py; do not edit
\begin{verbatim}
7 quantities that nothing in the artifact supports
   23,000   this is {} units, the price of 23,000 tokens of answer. The t
       26   the first formula gives 26 turns, and the script fi
    1.143   is {}, against $8/7 \approx 1.143$
5 quantities that repeat a generated value
    1,000   $T = 1,000$
       16   , $16/15$
32 quantities that repeat a literal of the source
   10,000    numbers.} Throughout we use one example: a system prompt of $h = 10,000
     0.08   tokens, and the prices $\rho = 0.08$
      200   kes from the price list of a provider. A plain loop of 200 turns then co
\end{verbatim}

\Description{An excerpt of the report of the tool for one paper: numerals with the text
around them, in three classes.}
\caption{A part of the tool's report for the latest version of paper 6: a few entries of each
class, each with the text in which the numeral stands. The figure is written by the script.}
\label{fig:excerpt}
\end{figure}

\paragraph{The report.} For every version the tool prints the quantities of each class with
the text around them, as in Figure~\ref{fig:excerpt}. The lists are short enough to read. An
unsupported quantity is something for the author to bind or to justify, and a repeated one
is a candidate for a name.

\paragraph{What the tool does not see.} Quantities written as words, such as ``twenty'' or
``a third'', are counted and not classified. Whole numbers up to \exSmallLimit{} are assumed
to be counts. And the tool does not know what a number means.

\section{Nine Papers}
\label{sec:audit}

\subsection{The Corpus}

The repository records every version of every paper with the commit that contains it. When
this paper was written there were \exVersions{} versions of \exPapers{} papers. All were
written by the same model in the same harness, over two days, and the later ones under a
stricter standard than the earlier ones: the repository began to ask for longer papers with
an artifact that tests every claim. The tool reads each version from the history, with the
artifact as it was at that commit.

\subsection{The Latest Versions}

\begin{table}[t]
\caption{The quantities in the prose of the latest version of each paper, by class, and the
share that is bound. The share in the first recorded version is given for comparison; a dash
means that the version has no quantities. The papers are numbered as in the repository.}
\label{tab:latest}
\small
\begin{tabular}{@{}lrrrrrrr@{}}
\toprule
Paper & version & bound & repeated & parameter & unsupported & \% bound & \% at first \\
\midrule
1 & \exVersiona & \exBounda & \exRepeata & \exParama & \exUnsupa & \exSharea & \exFirsta \\
2 & \exVersionb & \exBoundb & \exRepeatb & \exParamb & \exUnsupb & \exShareb & \exFirstb \\
3 & \exVersionc & \exBoundc & \exRepeatc & \exParamc & \exUnsupc & \exSharec & \exFirstc \\
4 & \exVersiond & \exBoundd & \exRepeatd & \exParamd & \exUnsupd & \exShared & \exFirstd \\
5 & \exVersione & \exBounde & \exRepeate & \exParame & \exUnsupe & \exSharee & \exFirste \\
6 & \exVersionf & \exBoundf & \exRepeatf & \exParamf & \exUnsupf & \exSharef & \exFirstf \\
7 & \exVersiong & \exBoundg & \exRepeatg & \exParamg & \exUnsupg & \exShareg & \exFirstg \\
8 & \exVersionh & \exBoundh & \exRepeath & \exParamh & \exUnsuph & \exShareh & \exFirsth \\
9 & \exVersioni & \exBoundi & \exRepeati & \exParami & \exUnsupi & \exSharei & \exFirsti \\
\bottomrule
\end{tabular}
\end{table}

Table~\ref{tab:latest} shows the latest version of each paper. Together they have
\exAllTotal{} quantities in their prose. Of these, \exAllBound{} are bound, which is
\exAllShare{} percent; \claim{three-quarters}{more than three quarters of the numbers that
a reader meets come out of an artifact by name}. \exAllRepeat{} repeat a value that the
artifact generates or records, \exAllParam{} repeat a literal of its source, and
\exAllUnsup{} are unsupported. Besides the quantities, the versions contain \exAllSmall{}
small whole numbers, \exAllFormal{} numerals inside formal objects, and \exAllWords{}
quantities written as words, which the tool does not classify.

\claim{six-least}{Among the papers that bind anything, the paper on the cost of an agent
loop binds the smallest share.} It works throughout with one example whose parameters, the
size of a prompt and the prices, are stated again wherever they are used, and each statement
is a numeral that repeats a literal of the script. They are the kind of number that goes
wrong silently: the day the example changes in the script, \exParamf{} places in the text
have to change with it.

Papers 2 and 5 bind nothing. Both are short papers without a script of their own, and both
have since been superseded by papers that have one.

\subsection{How the Numbers Got Bound}

\begin{table}[t]
\caption{The share of bound quantities, in percent, in every recorded version of every paper.
A dash means that the version has no quantities; an empty cell, that there is no such
version.}
\label{tab:history}
\small
\begin{tabular}{@{}lrrrr@{}}
\toprule
Paper & version 1 & version 2 & version 3 & version 4 \\
\midrule
1 & \exHistaa & \exHistab & \exHistac & \exHistad \\
2 & \exHistba & \exHistbb & \exHistbc & \\
3 & \exHistca & \exHistcb & \exHistcc & \\
4 & \exHistda & \exHistdb & \exHistdc & \exHistdd \\
5 & \exHistea & \exHisteb & \exHistec & \\
6 & \exHistfa & \exHistfb & & \\
7 & \exHistga & & & \\
8 & \exHistha & & & \\
9 & \exHistia & & & \\
\bottomrule
\end{tabular}
\end{table}

Table~\ref{tab:history} follows every paper through its versions. Two things changed the
numbers. Papers 1, 3 and 4 were rewritten when the repository raised its standard, and the
rewrite of each came with a script that writes its results to a generated file. And the
papers from the sixth on were written that way from their first version. A version that
changes nothing but the record of a paper, such as its keywords, keeps its share.

The last column of Table~\ref{tab:latest} shows where each paper started. The first paper is the
instructive case. Its first version reports a measurement of a library and contains
\exOneHand{} quantities, none of them bound. The artifact of that version wrote its results
to a file of text, and the author copied them into the tables. The tool finds \exOneCopied{}
of the \exOneHand{} numbers in that file, \exOneSource{} in the source, and \exOneDerived{}
nowhere; \claim{copied-most}{more than nine in ten were copies of the output}. The
others were ratios that the output printed in another unit, a quotient that the author
formed, and a size. We compared them with the output and the files by hand and found them
right.

So the first version of the first paper had no wrong number in it. It had \exOneHand{}
numbers that would have stayed as they were when the measurement was run again, and it was
run again: the second version reports a different sample, and only \exOneSame{} of the
\exOneHand{} numbers occur among its values. That
version binds them. The copying had been done without an error, once, and what the binding
bought was that it never had to be done again.

\subsection{The Unsupported Numbers}

We looked at each of the \exAllUnsup{} unsupported quantities of the latest versions.
\begin{itemize}
\item \exKindSize{} are sizes of an artifact: a paper says that its script or its
  mechanization has about so many lines.
\item \exKindDerived{} are derived by hand: the value of a formula at the example's
  parameters, a quotient of two bound costs, a fraction written as a decimal, a price
  converted into another unit.
\item \exKindCited{} are quoted from another paper of the repository, with its version.
\item \exKindExternal{} are prices from the price list of a provider, with a citation.
\item \exKindNotation{} are a parameter that the script writes as a fraction and the text as
  a decimal.
\item \exKindVersion{} is the version of a tool.
\end{itemize}
\claim{sizes-most}{The most common kind is the size of an artifact}, a number that nobody
thinks of as a result and that changes with every edit of a script.

\begin{table}[t]
\caption{The unsupported quantities of the latest versions, and what the script of this
paper compares each with. The two entries for paper 3 are the same parameter, stated twice.}
\label{tab:hand}
\small
\begin{tabular}{@{}lrll@{}}
\toprule
Paper & stated & what it is & compared with \\
\midrule
1 & \exHandSizea & lines of the mechanization & \exHandSizeA{} lines in the two files \\
2 & \exHandSizeb & lines of the mechanization & \exHandSizeB{} lines in the file \\
3 & \exHandSizec & lines of the script & \exHandSizeC{} lines \\
6 & \exHandSizef & lines of the script & \exHandSizeF{} lines \\
7 & \exHandSizeg & lines of the script & \exHandSizeG{} lines \\
2 & \exHandCitedbx & a median from paper 1 & a generated value of the cited version \\
2 & \exHandCitedby & a median from paper 1 & a generated value of the cited version \\
6 & \exHandDerivedfp & a cost as a number of tokens & \exHandQuestion{} units at four units a token \\
6 & \exHandDerivedfq & a formula at the example & \exHandFormula{} by the formula \\
6 & \exHandDerivedfr & eight sevenths as a decimal & \exHandFraction \\
6 & \exHandDerivedfs & a bound on a quotient of costs & the quotient, \exEgDerivedTrue \\
3 & \exHandNotationc & a rate of false acceptance & one in twenty in the script \\
9 & \exHandVersioni & the version of Lean & the header of the extracted graph \\
6 & \exHandExternalfu & the price of a cache write & not checked; cited \\
6 & \exHandExternalfv & the price of a cache read & not checked; cited \\
\bottomrule
\end{tabular}
\end{table}

For \exKindChecked{} of the \exAllUnsup{} the artifact of this paper now contains a check;
Table~\ref{tab:hand} lists them. It counts the lines of the files at the commit in question
and compares them with the stated size, within three percent. It evaluates the formulas. It
looks the quoted numbers up in the version that is cited. All of them pass. The two prices are not checked; they rest on
their citation. The checks are in the script of \emph{this} paper, which is the wrong place
for them: a later version of an audited paper is free to break them. What the audit shows is
that all of these numbers could have been bound, at the cost of a line each.

\subsection{Numbers That Point the Other Way}
\label{sec:refs}

An artifact also refers to its paper. A comment in a script says that a function tests a
theorem and gives the theorem's number, and the record of a paper says by number which
theorems a command checks. These numbers are typed, and nothing binds them: the numbering of
a paper is made by LaTeX, which the script does not see.

The tool collects every reference by number to a theorem, a lemma, a corollary, a
proposition, a definition, a table, or a figure in the scripts, the mechanizations, the
recorded output, and the record of a version. It then looks in the committed PDF of that
version for a statement, a table, or a figure of this kind with this number. The
\exVersions{} versions contain \exRefAll{} such references, and \exRefStale{} of them have
no target. \exRefStaleOne{} are in the fourth version of paper 1: \exRefStaleOneKinds{}
different references in the comments of its mechanization. \exRefStaleThree{} is in the
third version of paper 3. \claim{stale-rewritten}{Both are versions in which a paper was
rewritten.} Sections were added, the numbers of the theorems moved, and the comments kept the
old ones; one file of the mechanization had been copied from another paper together with
that paper's number for its theorem.

The tool finds only references without a target. A reference whose number has passed to
another statement goes through. When we corrected the two files of paper 1 we read their
comments, and found \exRefWrong{} more references of that kind: a comment that names the
theorem on merging by the number that the theorem on interfaces has now, and two that name a
proposition of the copied file by a number that belongs to another proposition. A check by
number cannot find these. A reference by label, or by the name of the theorem, could.

These are the only errors that the audit found. They mislead nobody about a result. But a
reader who follows the comment of a proof to the paper arrives at a different theorem, and
for a mechanization the comment is the only thing that says which statement of the paper a
formal statement is meant to be. We have not corrected the audited versions, which are what
they are; the corrections belong in new versions of the two papers.

The remedy is the same as for quantities, in the other direction. A build that writes the
numbers of the theorems to a file which the script reads, or a script that refers to
theorems by their labels and is checked against the list of labels, turns these references
into bound ones.

\subsection{What Another Seed Does}
\label{sec:seeds}

A bound number follows its artifact. How far would the numbers of these papers move if the
artifact were run on another sample? Six of the papers have a script that draws random
programs from one seed. We ran each script with \exSeedCount{} other seeds, at the commit of
its latest version, which makes \exSeedRuns{} runs. A run executes every test of the
theorems of its paper, so this is also a test of the tests.

\begin{table}[t]
\caption{The scripts of six papers, run with \exSeedCount{} other seeds: the number of
generated values, the largest number of them that changed in a run, and the number of runs in
which every test and every assertion of the script passed.}
\label{tab:seeds}
\small
\begin{tabular}{@{}lrrr@{}}
\toprule
Paper & generated values & change at most & runs that pass \\
\midrule
3 & \exSeedValuesc & \exSeedChangec & \exSeedPassc \\
4 & \exSeedValuesd & \exSeedChanged & \exSeedPassd \\
6 & \exSeedValuesf & \exSeedChangef & \exSeedPassf \\
7 & \exSeedValuesg & \exSeedChangeg & \exSeedPassg \\
8 & \exSeedValuesh & \exSeedChangeh & \exSeedPassh \\
9 & \exSeedValuesi & \exSeedChangei & \exSeedPassi \\
\bottomrule
\end{tabular}
\end{table}

Table~\ref{tab:seeds} shows the result. Of the \exSeedRuns{} runs, \exSeedPassed{} pass. No
test of a theorem fails in any run. In the papers that report counts over random programs,
many of the generated values change with the seed, as they should;
\claim{closed-none}{in the paper whose numbers are closed forms and exact searches, none
does}.

\claim{one-stops}{One run does not pass}, and it is the interesting one. It is a run of the script of
paper~8, and it stops at a line that asserts which of the weakened typing rules took the
most programs to refute. The paper says so in a sentence, and the author of the script had
added the assertion with a comment that the paper depends on it. With another seed a
different rule is the slowest. The sentence would then be false, and the script refuses to
produce the numbers. This is the claim of the calculus, written by hand in an earlier paper
before the calculus existed: a sentence about the numbers, with a test that fails the build
when the sentence stops being true. The same script has no such line for other sentences of
its paper, and for those nothing would have stopped.

The runs also say which numerals a typed copy would have survived. Take the generated
values of the scripts, with the seeds under which they pass, and call a numeral coarse if it
has at most two significant digits. We expected coarse numerals to survive a new sample more
often than precise ones. \claim{coarse-no-better}{They do not.} Of the \exSeedCoarse{} coarse
numerals, \exSeedCoarseKept{} are the same under every seed, which is \exSeedCoarsePercent{}
percent; of the \exSeedFine{} numerals with more digits it is \exSeedFineKept{}, or
\exSeedFinePercent{} percent. We read this as follows. What decides whether a number stays
is its origin and not its precision: a parameter and a closed form stay with all their
digits, and a number that was counted over a random sample moves, also when it is printed as
a percentage of two digits. For a typed number this means that rounding is no protection.

\subsection{The PDFs}

For every version the tool extracts the text of the committed PDF and looks for every
generated value that the source of that version uses. This needs a program that reads PDFs.
The \exVersions{} versions use \exPdfValues{} such values, and
\claim{pdf-agree}{every one of them appears in its PDF}. No PDF in the history is older than
its numbers. The check is one-sided: it shows that the right values are present, and not that
no other numbers are.

\section{Limits}
\label{sec:limits}

\paragraph{Meaning.} The audit is about where a numeral comes from, and says nothing about
what it measures. A script can compute a number carefully and the sentence around it can
describe another quantity. The calculus moves the reviewer's work to the right place and
does not remove it: for a paper that binds its numbers, what remains to be read is one line
of a script for each key. An earlier paper called the analogous set for a library of proofs
its audit surface~\citep{tpp1}.

\paragraph{Words.} The versions contain \exAllWords{} quantities in words. Some are the most
important sentences of their papers. The mechanism of claims covers them, but only where an
author uses it, and no tool will decide for an author that ``almost never'' needed a test.

\paragraph{Small numbers.} We treat whole numbers up to \exSmallLimit{} as counts of things
in the paper. Some of them are results. A paper that reports that a broken rule was refuted
after two tries has a result that the tool does not see unless it is bound, and in the
audited papers such numbers are bound.

\paragraph{Coincidence.} The classes \emph{repeated} and \emph{parameter} rest on equal
values. A numeral that equals some literal of a long script by chance is counted as
supported. The counts of unsupported quantities are therefore lower bounds.

\paragraph{One author.} The corpus is the work of one model in one harness over two days.
It shows what this model does under these instructions. It does not show what machine-written
papers are like in general. The audit found no wrong number, and that is a statement about a
small corpus written with artifacts in mind from the second day on.

\paragraph{Deceit.} The discipline constrains an author who wants to be constrained. It is
not a defense against one who does not.

\section{Related Work}
\label{sec:related}

\paragraph{Documents that compute.} Literate programming~\citep{knuth1984} puts a program
and its explanation in one source. Sweave~\citep{leisch2002} and its successors turn this
around for data analysis: the document contains code, and the numbers in the text are
results of running it. The case for publishing the code with the paper was made early by
\citet{buckheit1995}, and artifact evaluation made it an institution in our
field~\citep{krishnamurthi2015}. These systems give an author the means to bind a number.
They do not say which numbers of a text are bound and which are not, and that question is
what we formalize and measure.

\paragraph{Provenance.} Language-based provenance explains a part of an output by the parts
of the program and the input that it depends on~\citep{perera2012}, and has been used to link
the elements of a chart to the data behind them~\citep{perera2022}. Our labels are a coarse
provenance with two sources, the author and the artifact, and
Theorem~\ref{thm:provenance} is the corresponding noninterference
statement~\citep{volpano1996}. The finer question, which lines of a script a given number
depends on, is the reviewer's residue of Section~\ref{sec:limits}, and dependency analysis
is the tool for it.

\paragraph{Proofs in documents.} Proof assistants have long let a document quote a theorem
by name, so that the text shows the statement that was checked. Systems that interleave
prose with recorded output of the prover serve the same purpose for proofs that we pursue
for numbers~\citep{pitclaudel2020}. A theorem is easier than a number in one respect: its
statement is all there is. A number comes with rounding, which is why the calculus needs
intervals.

\paragraph{Carrying evidence.} The name follows proof-carrying code~\citep{necula1997},
where a program comes with evidence that a host checks before running it. An
evidence-carrying paper comes with the means to check its numbers, and the check is cheap in
the same way: it does not need trust in the producer.

\paragraph{Build systems.} That a committed result is still the result of its inputs is the
correctness condition of a build system~\citep{mokhov2018}. The comparison of recorded with
fresh definitions is a build check done by hand, and the gap between a PDF and its source is
a missing edge in that build.

\paragraph{This repository.} The audited papers are the first nine of the repository: two on
what a reader of a library of proofs must read~\citep{tpp1,tpp2}, and seven on programs that
call language models~\citep{tpp3,tpp4,tpp5,tpp6}, \citep{tpp7,tpp8,tpp9}. Their subjects do
not matter here. What matters is that each has a history of versions with artifacts, which makes an
audit of this kind possible after the fact.

\section{Conclusion}

A number in a paper is a claim that some computation gives this value. A reader should be
able to tell which claims the build has checked and which the author vouches for. For papers
written by a model the difference matters more than it did, because the author's care is the
part that does not scale, and binding is cheap. In nine papers by one model we found the
numbers right and about a fifth of them unbound. The unbound ones are parameters said
twice, sizes of scripts, and arithmetic done once. Each could be bound with a line. This
paper is the first of the repository in which each of them is.

\bibliographystyle{ACM-Reference-Format}
\bibliography{refs}

\end{document}
